% ============================================================
% ANEXO A - INSTRUMENTO DE AUDITORIA DE CALIDAD DEL ERS
% (Robinson Jeanpierre / jean200525 - PE5, Paso 1)
% Seis metricas derivadas de ISO/IEC 25010:2023 con formula,
% conteos, aritmetica visible, veredicto y accion de mejora.
% ============================================================
\section*{Anexo A. Instrumento de auditor\'ia de calidad del ERS}
\addcontentsline{toc}{section}{Anexo A. Instrumento de auditor\'ia de calidad del ERS}

\subsection*{A.1 Dise\~no del instrumento}

El instrumento aplica las seis m\'etricas operativas derivadas de ISO/IEC 25010:2023 definidas en la gu\'ia de la pr\'actica. Dos reglas lo gobiernan: toda m\'etrica se expresa como fracci\'on sobre un total contable (``si no se puede contar, no se puede reportar''), y los conteos base se publicaron en el repositorio \emph{antes} de calcular cualquier valor (\texttt{conteos\_base\_auditoria.csv} y \texttt{conteos\_base\_auditoria\_actualizado.csv}) \cite{iso25010,wohlin2012}. Cada m\'etrica se valora en una escala de 0 a 4:

\begin{itemize}[leftmargin=0.9cm]
  \item[4] Cumple el valor de referencia con evidencia completa.
  \item[3] Lo cumple con evidencia parcial.
  \item[2] Queda por debajo, pero con acci\'on de mejora documentada y aplicada.
  \item[1] Queda por debajo sin acci\'on de mejora.
  \item[0] No se midi\'o.
\end{itemize}

\noindent Pesos sugeridos por la gu\'ia: M1 $=$ 0,20; M2 $=$ 0,15; M3 $=$ 0,20; M4 $=$ 0,20; M5 $=$ 0,10; M6 $=$ 0,15.

\subsection*{A.2 Conteos base publicados}

\begin{center}
\begin{tabular}{@{}p{6.8cm}@{\;}p{3.4cm}@{\;}p{3.4cm}@{}}
\hline
\textbf{Conteo} & \textbf{L\'inea base auditada} & \textbf{Tras fusi\'on IA-02} \\
\hline
Requisitos funcionales & 27 & 31 \\
Requisitos no funcionales & 15 & 26 \\
Total de requisitos & 42 & 57 \\
Casos de uso especificados & 16 & +CU de RF-28..31 en cierre \\
Actores operativos con RF & 2 & 2 \\
Historias de usuario en el ERS & 16 & 16+ \\
Pares intra-\'ambito medidos para M2 & 116 (ocho \'ambitos) & 191 (nueve \'ambitos) \\
Pares te\'oricos $n(n-1)/2$ (referencia) & 861 & 1596 \\
Filas vigentes en matriz al medir & 28/42 & cierre Paso 2 \\
\hline
\end{tabular}
\end{center}

\subsection*{A.3 Fichas por m\'etrica}

\subsubsection*{M1 --- Completitud}
\textbf{F\'ormula:} $M1a = (\#\,\text{req.\ con 4 atributos}/\#\,\text{req.\ totales})\times100$; $M1b = (\#\,\text{CU espec.}/\#\,\text{CU ident.})\times100$; $M1c = (\#\,\text{actores con}\geq1\,\text{RF}/\#\,\text{actores})\times100$.\\
\textbf{Aritm\'etica:} $M1a = 42/42\times100 = 100\,\%$ en la l\'inea base; la primera medici\'on sobre la versi\'on 2.0 dio $54/57\times100 = 94{,}7\,\%$ (Origen sin fuente trazable en RNF-20, RNF-24 y RNF-25); tras a\~nadir el requisito de origen a esas tres fichas, $57/57\times100 = 100\,\%$. $M1b = 16/16\times100 = 100\,\%$; $M1c = 2/2\times100 = 100\,\%$.\\
\textbf{Referencia:} $M1a\geq95\,\%$; $M1b=M1c=100\,\%$. \textbf{Veredicto:} cumple. \textbf{Nivel:} 4. \textbf{Acci\'on aplicada:} correcci\'on de origen en tres fichas de equidad y re-medici\'on, con el par $94{,}7\,\%\rightarrow100\,\%$ publicado en \texttt{aritmetica\_metricas.csv}.

\subsubsection*{M2 --- Consistencia}
\textbf{F\'ormula:} $M2 = 1 - \#\,\text{conflictos}/\#\,\text{pares analizados}$.\\
\textbf{Aritm\'etica:} ola 1: 116 pares intra-\'ambito generados por script (ocho \'ambitos, $n=42$), cero conflictos $\Rightarrow M2 = 1 - 0/116 = 1{,}00$. Ola 2 sobre la versi\'on 2.0: 191 pares intra-\'ambito (nueve \'ambitos, $n=57$), cero conflictos $\Rightarrow M2 = 1 - 0/191 = 1{,}00$. El conteo te\'orico todos-contra-todos ($n(n-1)/2$: 861 y 1596) queda registrado solo como referencia en \texttt{conteos\_base\_auditoria\_actualizado.csv}.\\
\textbf{Referencia:} $\geq0{,}98$. \textbf{Veredicto:} cumple. \textbf{Nivel:} 4. \textbf{Acci\'on:} observaciones de solapamiento documentadas por par en \texttt{m2\_pares\_analizados.csv}; DR-02 fue hallado y cerrado en la re-inspecci\'on previa a ambas olas de medici\'on.

\subsubsection*{M3 --- Verificabilidad}
\textbf{F\'ormula:} $M3 = (\#\,\text{requisitos con criterio comprobable}/\#\,\text{totales})\times100$.\\
\textbf{Aritm\'etica:} $57/57\times100 = 100\,\%$ sobre la versi\'on fusionada ($42/42$ en la l\'inea base): umbral, unidad y m\'etodo presentes en todo el cat\'alogo, incluidas las 15 fichas de IA incorporadas en PE5.\\
\textbf{Referencia:} $\geq90\,\%$. \textbf{Veredicto:} cumple. \textbf{Nivel:} 4. \textbf{Acci\'on:} sin acci\'on; regla ``todo requisito nace con su m\'etrica''.

\subsubsection*{M4 --- Trazabilidad}
\textbf{F\'ormula:} $M4ade = (\#\,\text{RF con cadena adelante completa}/\#\,\text{RF})\times100$; $M4atr = (\#\,\text{RF con fuente identificada}/\#\,\text{RF})\times100$.\\
\textbf{Aritm\'etica:} $M4atr = 57/57 = 100\,\%$ en la versi\'on 2.0 ($42/42$ en la l\'inea base); $M4ade = 16/16 = 100\,\%$ sobre Must Have (RF-28 a RF-31 son Deber\'ia-tener y quedan fuera del alcance). La primera medici\'on de cobertura global de filas vigentes arroj\'o $28/42 = 66{,}7\,\%$, bajo la referencia, por filas de RF-17 a RF-27 pendientes de reconstrucci\'on con las diez columnas.\\
\textbf{Referencia:} $M4ade\geq90\,\%$; $M4atr=100\,\%$. \textbf{Veredicto:} subm\'etricas cumplen; cobertura global en cierre por Paso~2 (Danela Arteaga). \textbf{Nivel:} 2 (debajo en cobertura global, acci\'on documentada y en ejecuci\'on).

\subsubsection*{M5 --- Modificabilidad}
\textbf{F\'ormula:} acoplamiento promedio $=$ requisitos afectados al modificar cada uno de una muestra de $\geq5$ requisitos representativos, recorriendo la matriz.\\
\textbf{Aritm\'etica:} muestra de 5 requisitos con promedio inicial $3{,}50$; tras re-especificar las vistas derivadas (referenciar el requisito base en vez de duplicarlo), re-medici\'n $= 3{,}00$ (detalle por requisito en \texttt{calculo\_M5\_modificabilidad.csv}).\\
\textbf{Referencia:} $\leq3{,}0$. \textbf{Veredicto:} cumple tras correcci\'on. \textbf{Nivel:} 4. \textbf{Acci\'on aplicada:} eliminaci\'on de duplicaciones; par antes/despu\'es reportado en Secci\'on~8.

\subsubsection*{M6 --- Correcci\'on}
\textbf{F\'ormula:} $M6 = \#\,\text{defectos hallados despu\'es de la inspecci\'on}/\#\,\text{requisitos totales}$.\\
\textbf{Aritm\'etica:} re-inspecci\'on inicial sobre el ERS corregido en PE4: $5/42 = 0{,}119$ defectos/requisito; tras cerrar DR-01 a DR-05, conteo residual $1/42 = 0{,}024$.\\
\textbf{Referencia:} $\leq0{,}05$. \textbf{Veredicto:} cumple tras correcci\'on. \textbf{Nivel:} 4. \textbf{Acci\'on aplicada:} correcciones DR-02 a DR-04 documentadas con commit de evidencia.

\subsection*{A.4 Tabla consolidada del instrumento}

\begin{center}
\begin{tabular}{@{}l l l l l l l@{}}
\hline
\textbf{M\'etrica} & \textbf{Valor} & \textbf{Ref.} & \textbf{Cumple} & \textbf{Nivel} & \textbf{Peso} & \textbf{Aporte}\\
\hline
M1 Completitud & 100\,\% & $\geq95$\,\% & S\'i & 4 & 0,20 & 0,80 \\
M2 Consistencia & 1{,}00 & $\geq0{,}98$ & S\'i & 4 & 0,15 & 0,60 \\
M3 Verificabilidad & 100\,\% & $\geq90$\,\% & S\'i & 4 & 0,20 & 0,80 \\
M4 Trazabilidad & ade 100\,\%; atr 100\,\%; global en cierre & $\geq90$\,\% & Parcial & 2 & 0,20 & 0,40 \\
M5 Modificabilidad & 3,50 $\rightarrow$ 3,00 & $\leq3{,}0$ & Tras corr. & 4 & 0,10 & 0,40 \\
M6 Correcci\'on & 0,119 $\rightarrow$ 0,024 & $\leq0{,}05$ & Tras corr. & 4 & 0,15 & 0,60 \\
\hline
\multicolumn{6}{@{}r}{\textbf{Resultado ponderado del instrumento}} & \textbf{3,60/4,00} \\
\hline
\end{tabular}
\end{center}

\noindent Resultado ponderado: $(0{,}80+0{,}60+0{,}80+0{,}40+0{,}40+0{,}60) = 3{,}60$ sobre 4,00 ($90\,\%$). El \'unico nivel no m\'aximo es M4, cuya cobertura global depende del cierre del Paso~2; las subm\'etricas exigibles ya cumplen su referencia.

\subsection*{A.5 Reproducibilidad de la auditor\'ia}

Los conteos son regenerables desde la ra\'iz del repositorio con los scripts \texttt{conteo\_requisitos.ps1} (acepta el par\'ametro \texttt{-ErsPath} para apuntar a la versi\'on a auditar; por omisi\'on mide la l\'inea base v1.0) y \texttt{generar\_pares\_m2.ps1} (PowerShell 5.1+); sus salidas en \texttt{analisis\_estadistico/salidas/} son deterministas para una misma versi\'on del ERS, verificado por comparaci\'on de hash SHA-256 (\texttt{verificacion\_reproducibilidad.md}). Los insumos de M5 y M6 viven en \texttt{calculo\_M5\_modificabilidad.csv}, \texttt{registro\_defectos\_lineabase\_1B\_2A.csv} y \texttt{registro\_reinspeccion\_PE5.csv}. Cualquier evaluador puede replicar cada cifra de este anexo sin acceso a los autores \cite{wohlin2012}.
